\begin{abstract}
Model-based data assimilation (DA) often loses to estimators that use less
physics. Observation-only products can beat ocean reanalyses, learned priors can
beat the true dynamical model, and neural schemes are routinely reported to beat
classical ones. The usual explanations --- a wrong model, a lossy algorithm, a
poor tuning --- are seldom separated.
%
We treat every DA scheme, classical or learned, as an estimator of the
conditional distribution of the state given its context. We decompose its error
along a chain of idealisations. The \emph{irreducible} term is what no scheme can
beat. A \emph{restriction} is the error of using less information than is
available. A \emph{misspecification} is the error of believing a wrong model. An
\emph{approximation} is the error of computing one's own target inexactly.
%
We extend the chain from the mean to the reported uncertainty with a
flow-matching score, $\FMS_\tau$. It equals the mean-squared error at
$\tau = 0$, is strictly proper for $\tau \in (0, 1)$, and splits exactly into an
irreducible part, a restriction part and a remainder at every $\tau$.
%
On three case studies --- a two-scale Lorenz-96 system carrying the full
scheme matrix, a Lorenz-63 system for the model-error axis and a
quasi-geostrophic ocean for realism --- with a perfect model and with model
error, under matched information and tuning parity, three results stand out.
\begin{itemize}[leftmargin=*,nosep]
  \item With a perfect model, the deficit of classical DA is mostly restriction
    and approximation. Smoothing removes 16--19\,\% of the filter's RMSE, and
    going from 30 to 100 ensemble members removes 38\,\% of the smoother's
    mean-squared error.
  \item Under model error, misspecification dominates. It shows in the
    large-$\tau$ score before it shows in the RMSE: the smoother's calibration
    loss grows from 34\,\% to 63\,\%. The marginal value of observations for
    strong-constraint 4D-Var collapses sevenfold, and on the quasi-geostrophic
    ocean the best perfect-model scheme, strong-constraint 4D-Var, becomes the
    worst.
  \item A learned prior beats the true-model hard constraint and the 30-member
    smoother, but not the 100-member one. The ordering is set by approximation,
    not by the fidelity of the prior.
\end{itemize}
A restriction costs variance honestly and preserves calibration; a
misspecification costs bias silently and destroys it.
\end{abstract}
